\chapter{DISCUSSION AND LIMITATIONS}
\label{ch:discussion}
\section{What the Study Can Explain}
Follow My Voice asks whether a synthetic voice resembling the listener changes search efficiency or experience relative to a neutral voice the listener has already chosen. Because participant results are pending, each interpretation is stated conditionally, as a reading the data could support rather than a conclusion about which voice is better.

The discussion distinguishes three kinds of evidence. The primary question concerns measured efficiency, the secondary questions concern workload and motivation, and the exploratory questions examine participants' evaluations of the voices and their reported reliance on guidance. Agreement among these outcomes could strengthen a design interpretation, whereas disagreement could reveal a tradeoff that a single outcome would conceal.

\section{Interpreting Efficiency and Experience Together}
If completion time favors the self-similar voice, interpretation considers the estimate's uncertainty together with the manipulation check and guidance quality. A difference accompanied by systematic intelligibility or delivery problems in the neutral condition could reflect those problems rather than anything specific to self-similarity. If both conditions have comparable delivery quality and participants recognize the clone, the result would more directly support a voice-condition effect in this task, although it would still not separate resemblance from other differences between the chosen and synthesized voices.

Faster completion with the neutral voice would be consistent with a chosen external voice being easier to follow, but also with intelligibility or delivery problems in the cloned condition. Delivery records and acoustic-quality reporting could help distinguish these readings, although the study could not establish the mechanism without further evidence. Liking, naturalness, and reliance ratings may help characterize the experience in either case, but a correlation between such ratings and performance would not show that they mediate the voice effect.

If completion times are similar, the width of the uncertainty interval becomes central. A wide interval would leave practically important differences unresolved, whereas a narrow interval near the null could constrain the size of a plausible effect under the tested conditions. A nonsignificant test alone does not establish equivalence. Equivalence testing is not part of the current analysis specification, and adding it would require a meaningful bound and procedure defined before inspecting the result.

A preference for self-similar speech may occur without a detectable improvement in completion time, while efficient search may coexist with discomfort.

Kao et al. examine an avatar voice in a setting where continued voluntary play can reflect engagement, whereas Guo et al. report favorable impressions alongside eeriness in virtual collaboration~\cite{kao2021selfsimilar,guo2024doppelganger}. Here, a voice instructs the listener through a fixed search task. Faster completion therefore addresses efficiency, while enjoyment, discomfort, and debrief responses address the experience of following it.

\section{Practical Relevance and Transfer to Other Search Tasks}
The applications motivating the introduction identify where guided search could be useful, rather than domains already validated by this experiment. Locating task-relevant items in maintenance or training, finding a book, and noticing relevant VR content share a need to connect assistance with something in the environment~\cite{binetti2021outofview,zhang2023seehear,lange2020hivefive}. The present study examines a voice-choice question within that interaction; it does not test industrial task completion, learning transfer, navigation safety, or narrative comprehension.

The practical value of the study depends on how its outcomes combine. Faster search accompanied by acceptable workload and a comfortable voice experience would support considering self-similar guidance for comparable tasks. A favorable experience without a detected speed difference could still motivate offering the clone as an optional voice, though not claiming that it makes search more efficient. Discomfort, intelligibility problems, or a preference for the neutral voice would instead mark limits to that option.

Transfer also depends on how much information users already have. Here, the participant must use spatial guidance to infer an initially unnamed target; in Zhang et al.'s book task, the target title was supplied~\cite{zhang2023seehear}. A known-target task may allow people to continue independently and encounter the voice less often. Real work can additionally require recognizing meaningful parts, remembering procedural steps, moving through space, or coordinating with another person. Replication in those contexts would need to establish both the quality of guidance and the effect of voice choice.

\section{Voice Identity and Acoustic Confounds}
The voice conditions may differ in acoustic properties as well as perceived resemblance. The self-similar voice depends on recording quality and synthesis, whereas the neutral voice comes from a heterogeneous catalog, and the two may differ in pitch, prosody, speaking rate, pronunciation, and clip duration. Because the same words and provider do not guarantee matched acoustics, the study documents these properties and pilot intelligibility at comfortable playback levels rather than claim exact equivalence.

The variety of chosen neutral voices is a deliberate consequence of the design: requiring every participant to hear the same neutral voice would replace the comparison with a preferred alternative by a comparison with a fixed one. The analysis should retain chosen voice identifiers and describe the distribution, while avoiding an overfitted model with unsupported estimates for sparsely chosen voices. An exploratory association between a neutral voice and performance may reflect who selected it as well as its acoustic properties.

The current self-similar manipulation applies no additional pitch shift, panning, or inner-thought treatment. Earlier development attempts to make cloned speech sound like an internal monologue are not part of the study, and reintroducing such an effect in only one condition would change the manipulation and require a revised rationale and analysis.

\section{Adaptive Guidance as a Shared Mechanism}
Gaze-dependent sonification, staged spatial cues, and audiovisual assistance investigate different ways to connect a cue with a search target~\cite{losing2014guiding,kwok2022twostep,binetti2021outofview}. Follow My Voice adapts these precedents into spoken orientation and refinement so that the assigned voice participates throughout the task. A difference between the two voice conditions would not demonstrate that speech outperforms those techniques, or that Jev improves search: neither comparison is part of the participant design.

A shared policy does not produce identical spoken sequences. Because participants turn, pause, look, and point differently, the policy supplies different directions and can emit different numbers of hints.

Nevertheless, policy drift would threaten the comparison. Changing phrases, timers, or provider behavior midway through collection could alter assistance quality. The study release must pin the application build, configuration, phrase semantics, voice preparation settings, and decision-provider version where possible. Any unavoidable service change must be documented and considered in the analysis.

The central state machine improves auditability by assigning guidance evidence and transitions to a single local component, but it cannot guarantee correct guidance. Calibration error can produce an incorrect observation, an unclear reference frame can produce an unhelpful phrase, and a stale response can pass an incomplete check.

\section{Practice, Learning, and Order}
One practice trial reduces setup burden but provides limited exposure to a complex interaction in which a participant must learn to hold gaze on the readiness cross, follow speech, use the chair comfortably, aim the right controller, and press the trigger to select. A fixed single trial does not prove that every participant has acquired these skills equally, and the practice-proficiency rule remains to be established through pilots.

Practice uses the first assigned voice, which aligns initial exposure with the first block but also couples practice familiarity with period. Counterbalancing helps distribute this coupling across voice conditions; it does not prevent fatigue, learning, novelty, or asymmetric carryover. These limitations are reported with period and order effects and any relevant sensitivity analysis.

\section{Retrospective Questionnaires and Burden}
The final-only survey avoids interrupting the two-block search sequence with a long questionnaire, but it also requires participants to recall Block 1 after completing Block 2. Ratings may therefore reflect recency, contrast, or imperfect recall. Each rating set identifies its referenced block; the reminder linking that block to its voice remains to be piloted. These reminders may aid recall but cannot guarantee accurate recollection.

The condition-specific battery contains 130 numeric responses across both blocks, in addition to setup ratings, comparison questions, and debrief. Its completion time, missingness, and comprehension therefore require pilot measurement. Because participants may be tired or uncomfortable after exposure, researchers check whether they are comfortable continuing before administering the retrospective sets.

Adapting an embodied-agent measure to a voice agent can alter its meaning. Co-presence or anthropomorphism items may be understood differently without a visible character. The thesis retains attribution and wording changes, reports the adaptation, and avoids asserting new psychometric validation from a small experimental sample. Custom reliance and naturalness items remain individual exploratory measures.

\section{Gaze Quality and Feedback Coupling}
Gaze serves both as an input to assistance and as a recorded behavioral signal. Tracking gaps can change the participant's experience and reduce measurement coverage, affecting both the outcome being measured and the evidence available to interpret it. Coverage and technical events are considered alongside performance without automatically being treated as pre-existing participant characteristics.

A highlighted vicinity candidate can help indicate what the system currently attends to, but it can also suggest a reference object that differs from the mathematical gaze ray. For example, an instruction to look lower can be consistent with a gaze ray above the target yet confusing if the target appears higher than a highlighted neighbor. A cue should be evaluated in its visible interaction context, not solely against the sign of a geometric delta. Development-only rings and diagnostic markers must not silently enter the participant condition.

The proposed validity-aware stream and independent calibration checks address these risks, but headset-specific spatial performance and the rate of distinct gaze observations remain empirical questions. Validation results for the Vive Pro Eye cannot substitute for Focus Vision measurements, and an interaction threshold chosen for usability does not establish a validated measure of cognition.

\section{Privacy, Consent, and Reproducibility}
Coding participant identifiers does not make voice recordings anonymous, because both natural and synthesized speech can remain recognizable. The consent and data-management documents must describe recording, synthesis, cloud processing, storage, access, and deletion using the services actually employed. Older local-only descriptions are inconsistent with current ElevenLabs voice preparation and Jev decisions and require reconciliation before participant use.

The thesis reports the relevant processing stages without including credentials or identifiable recordings. A reproducibility package can provide code, versioned schemas, example synthetic records, and de-identified summaries where authorized. Distribution of cloned voices or raw participant audio requires a separate basis and must not be assumed from consent to participate.

Cloud dependencies also affect replication. A future researcher may not obtain the same voice model or identical decision behavior. Preserving prepared-library metadata and actual playback records helps describe the original exposure, while a constrained phrase interface makes alternative providers testable; neither removes the need to report provider changes.

\section{Future Work on Self-Similar Assistance}
Future studies could examine repeated exposure, task context, and the separate contributions of perceived resemblance and acoustic naturalness. Manipulating resemblance while controlling other acoustic properties would address a narrower causal question than the present comparison with a participant-selected neutral voice.

Other extensions could examine a visible agent body, a known-target search task, or a broader participant population. Each would change how the voice is encountered and would require a revised rationale and method.

Memory, provider autonomy, and gaze-guidance policy could also be varied in separate systems studies. Those comparisons would concern the supporting mechanism and must not be inferred from the present voice contrast. Similarly, eye-event or pupil analyses should be added only when validated measurements answer a defined question about task demands or voice experience.

\subsection{Gaze-Informed Adaptation of Spoken Assistance}
Nerella et al.'s gaze-based prediction study also motivates testing whether momentary task demands can inform when an agent speaks~\cite{nerella2026gaze}. For Follow My Voice, such an extension requires task-specific labels, evaluation on unseen participants, and a separate test of adaptive hint timing. The cited work leaves real-time deployment evaluation open; its classification result does not establish that adapting spoken guidance improves search.

\subsection{Complementing Behavioral Measures with Brain Signals}
A future extension could combine electroencephalography (EEG) with gaze, task performance, and questionnaires to investigate the attentional demands of receiving self-similar and neutral spoken guidance. Yang et al.'s \emph{Synchronizing with Attentional Sampling} provides a relevant precedent: the authors use EEG in immersive VR to investigate attentional sampling during gaze and hand--eye interaction as part of a study of visual flicker guidance~\cite{yang2026attentional}. This work motivates examining neural activity alongside observable search behavior, although its visual-guidance findings do not establish how voice resemblance affects attention or how frequently spoken hints should occur.

For the present voice comparison, synchronized recordings could relate neural activity to actual instruction onset, orientation toward the target, and controller selection. A follow-up study could then examine whether the two voices are associated with different responses during these stages, even when completion times are similar. Such measurements would complement the participant's reports of workload and experience rather than serve as direct measures of enjoyment, trust, or perceived voice similarity.

This extension would first require a feasibility study of EEG and headset compatibility, event synchronization, and the effects of eye movements, head turns, and muscle activity on the recordings. Neural measures and analysis windows would need to be specified in advance, with comparable task demands and speech exposure across voice conditions.
